%% kapa: exam=1 round=33 year=2022 period=2 form=A subject=appraisal_law subject_ko=감정평가관계법규 seq=27 number=27
%% source: Q-Net 감정평가사 기출문제 (cst00309) artlSeq=5209997 file=raw/1008/5209997/2022년도 제33회 감정평가사 1차시험 2교시(감정평가 관계법규, 회계학).pdf
%% conversion: kordoc 4.24.51 -> tools/gen1.py (auto); review=proofread-sonnet; flags=none
\begin{question}{27}
건축법령상 공개공지등에 관한 설명으로 옳지 않은 것은? (단, 조례는 고려하지 않음)
\begin{choices}
  \item 공개공지등은 해당 지역의 환경을 쾌적하게 조성하기 위하여 일반이 사용할 수 있도록 설치하는 소규모 휴식시설 등의 공개 공지 또는 공개 공간을 지칭한다.
  \item 공개공지등은 상업지역에도 설치할 수 있다.
  \item 공개 공지는 필로티의 구조로 설치할 수 있다.
  \item 숙박시설로서 해당 용도로 쓰는 바닥면적의 합계가 3천 제곱미터인 건축물의 대지에는 공개 공지 또는 공개 공간을 설치하여야 한다.
  \item 판매시설 중 「농수산물 유통 및 가격안정에 관한 법률」에 따른 농수산물유통시설에는 공개공지등을 설치하지 않아도 된다.
\end{choices}
\end{question}
\answer{4}
